% Introduction de la synthèse : la feuille de route (v2/ROADMAP.md), écrite à la main.
\part{Feuille de route}

\section{Objectif}
\label{syn-objectif}
Ce document rassemble le travail de recherche du projet « jumeau numérique de l'offre de soins »
au 14 novembre 2026. Le projet est né dans un hackathon de 48 heures (10 personnes) et a été
recadré le même jour en projet de recherche quantitative.

\paragraph{La question.} Simuler des \textbf{politiques publiques nationales}, chacune exprimée
comme un réglage d'hyperparamètres sourcés, et mesurer leur effet à une date future $t$
(horizon de référence \textbf{2050}) sur l'effectif \enquote{libéral} de cinq professions :
\begin{itemize}
  \item les sages-femmes (maïeutique) ;
  \item les chirurgiens-dentistes (odontologie) ;
  \item les pharmaciens ;
  \item les masseurs-kinésithérapeutes ;
  \item les infirmiers.
\end{itemize}
La médecine générale est hors du périmètre. Les spécialités médicales (cardiologie,
dermatologie, ophtalmologie) forment une phase~2.

\paragraph{Définitions retenues.}
\begin{itemize}
  \item \textbf{« Libéral »} désigne le secteur privé, lucratif ou non. Tout employeur public
        est exclu : hôpital public, PMI, Éducation nationale, EHPAD publics, centres de santé
        municipaux. Les pharmaciens adjoints d'officine, les salariés de cliniques privées et
        l'industrie sont donc inclus. Ce n'est pas le sens des sources (RPPS, ameli), où
        « libéral » veut dire exercice indépendant ; chaque sortie rappelle la définition.
  \item \textbf{Échelles} : les modèles tournent à tous les niveaux, de la commune au national,
        parce que la régulation de l'installation s'applique par commune et que le directeur
        général d'ARS doit avoir la visibilité maximale.
  \item \textbf{Horizon} : 2050. Les données allant de 2012 à 2026, aucun backtest n'est
        possible au-delà d'environ 10 ans. La validation à 2050 passe donc par la
        reproduction des projections de référence (DREES) et par des limites théoriques
        (axe~3).
\end{itemize}

\paragraph{Principes non négociables.}
\begin{itemize}
  \item \textbf{Démarche scientifique.} Chaque méthode et chaque résultat est reproductible
        (commande, données, millésime, version du code) et falsifiable : le test qui peut
        échouer est écrit avant de le lancer.
  \item \textbf{Aucune boîte noire.} Chaque paramètre a un sens, une source, une unité et une
        incertitude, et un statisticien extérieur doit pouvoir vérifier un modèle à partir de
        son article seul.
  \item \textbf{Code en briques.} Chaque mécanisme est un module (plug-in) qu'on ajoute,
        retire ou remplace sans toucher aux autres.
  \item \textbf{Données ouvertes uniquement}, chacune avec son jeu de données et son
        millésime.
  \item \textbf{Lexique neutre} : on décrit et on compare (« accès moindre »), on ne juge pas.
  \item \textbf{Plusieurs estimateurs comparés}, jamais un chiffre unique présenté comme
        « la » projection.
  \item \textbf{Provenance de chaque nombre} : mesuré (\mesure), dérivé (\derive) ou supposé
        (\suppose).
\end{itemize}

\section{Modèles de référence et structure}
Le jumeau reprend les hyperparamètres de trois modèles publics :
\begin{itemize}
  \item \textbf{Omphale} (Insee) pour la population : fécondité, espérance de vie, solde
        migratoire ;
  \item \textbf{Hippocrate} (DREES) pour les médecins : départs, mortalité, âge à
        l'installation, entrées formées en France et diplômés étrangers ;
  \item le \textbf{Dossier DREES n°76} (2021), modèle par composantes avec transitions entre
        modes d'exercice pour les dentistes, sages-femmes et pharmaciens.
\end{itemize}
Chaque politique publique modifie un hyperparamètre nommé.

La structure commune est une identité comptable par cohorte : l'effectif d'une année est celui
de l'année précédente, plus les entrées, moins les sorties, plus les transitions.
\begin{itemize}
  \item \textbf{Ce qui est connu} : le vieillissement, les étudiants déjà inscrits, la
        population projetée.
  \item \textbf{Ce qui est estimé} : la rétention, les entrées, la répartition territoriale,
        le passage du privé au public.
\end{itemize}
L'axe~2 cherche de meilleurs estimateurs \emph{sur cette même structure transparente}, au lieu
de la remplacer.

\section{Six axes de recherche}
\begin{center}
\small
\begin{tabular}{p{0.15\linewidth}p{0.37\linewidth}p{0.38\linewidth}}
\toprule
Axe & Question & Livrable \\
\midrule
0. Jumeau & Quel socle commun pour brancher méthodes, paramètres et données ? & état, modules de flux, scénarios ; critères d'acceptation capables d'échouer \\
1. Hyperparamètres & Quelles spécificités par profession, sourcées et mesurables ? & fiches de paramètres ; catalogue de politiques simulables \\
2. Méthodes & Quels modèles prévoient mieux sans être des boîtes noires ? & estimateurs branchés sur le jumeau, comparés sur un banc \\
3. Validation & Comment montrer qu'un modèle et un levier sont justes, jusqu'en 2050 ? & protocoles pré-enregistrés et gelés (N1 à N6) \\
4. Données & Quelles séries et quels événements manquent ? & séparation public / privé, registre d'événements, zonage, grain communal \\
5. Sensibilité & Quels hyperparamètres déplacent vraiment la sortie, et avec quelle fourchette ? & indices de Morris et de Sobol, fourchettes 2050 \\
\bottomrule
\end{tabular}
\end{center}

Les axes~1, 3 et~4 produisent des documents et des données, et ont démarré en même temps que le
jumeau. Les axes~2 et~5 ont démarré une fois son contrat stabilisé.

\section{Comment le travail a été contrôlé}
\begin{itemize}
  \item \textbf{Revues adversariales du jumeau.} Trois revues indépendantes, chacune menée par
        9 agents en lecture seule, ont soumis chaque constat à un agent chargé de le réfuter en
        exécutant le code.
        \begin{itemize}
          \item Première revue : 6 défauts confirmés. Plusieurs critères annoncés étaient vrais
                par construction.
          \item Deuxième revue : 4 défauts bloquants, dont un levier retraite qui surestimait
                son effet de 41~\% à $\Delta = 2$ ans et une régulation qui créait des
                professionnels.
          \item Troisième revue : plus de défaut bloquant ; la conservation est vérifiée à
                0,5 personne près en 2050 (\mesure).
        \end{itemize}
        Ces revues sont aussi un cas d'étude pour l'axe~3 : des tests qui passent par
        construction.
  \item \textbf{Protocoles gelés avant calcul.} Les empreintes SHA-256 des protocoles de
        l'axe~3 sont calculées avant tout résultat.
  \item \textbf{Registre des décisions.} Chaque choix, fait, remplacé ou rejeté, est daté
        (\texttt{v2/docs/lab/decisions.md}), et chaque question ouverte a un coût
        (\texttt{v2/docs/lab/open.md}).
\end{itemize}

\section{Lire ce document}
Les parties suivantes reprennent \emph{tels quels} les six prétirages d'axe : seuls les renvois
internes et les clés bibliographiques sont préfixés. Chaque partie s'ouvre sur le résumé de son
axe. Viennent ensuite :
\begin{itemize}
  \item une partie sur l'\textbf{interface et les tests utilisateurs} ;
  \item une partie sur \textbf{ce qui reste à faire}.
\end{itemize}
Le document est régénéré par \texttt{v2/papers/synthese/build.py} à partir des articles
d'axe : il suit leurs mises à jour.
